% Chapter 4 -- Rubric: PO3 (3.3, 3.4), PO5 (5.1)
\chapter{System Design and Implementation}
\label{ch:design}

\section{Functional Design and Sub-problem Partitioning}
\label{sec:architecture}

The problem decomposes into five subsystems with a deliberate dependency
order. The knowledge ontology is upstream of everything: it defines the skills
the engine reasons about, the topics the interface displays, and the
specifications the content pipeline must generate against. The inference engine
depends on the ontology but not on the content; the content pipeline depends on
the ontology but not on the engine; the API composes engine and content; the
interface depends only on the API. This ordering is what allowed the ontology to
be rebuilt late in the project (\Cref{sec:ontology-rebuild}) without rewriting
the engine.

\begin{figure}[htbp]
  \centering
  \begin{tikzpicture}[
      font=\small,
      box/.style={draw, rounded corners=2pt, align=center, minimum height=9mm,
                  inner xsep=3mm, fill=gray!8},
      store/.style={draw, align=center, minimum height=9mm, inner xsep=3mm,
                    fill=gray!18},
      off/.style={draw, dashed, rounded corners=2pt, align=center,
                  minimum height=9mm, inner xsep=3mm},
      arr/.style={-{Stealth[length=2mm]}, thick}
    ]
    % --- serving path -------------------------------------------------
    \node[box] (ui) at (0,0) {React + Vite client\\{\scriptsize diagnostic, practice, mastery views}};
    \node[box] (api) at (0,-2.1) {FastAPI application\\{\scriptsize \texttt{server.py} --- 14 endpoints}};
    \node[box] (engine) at (0,-4.2) {Adaptive engine\\{\scriptsize BKT $+$ mastery policy $+$ diagnostic}};
    \node[store] (db) at (0,-6.3) {PostgreSQL (Supabase)\\{\scriptsize learners, mastery, questions, ontology}};

    \draw[arr] (ui) -- node[right=1mm]{\scriptsize HTTP / JSON} (api);
    \draw[arr] (api) -- node[right=1mm]{\scriptsize in-process calls} (engine);
    \draw[arr] (engine) -- node[right=1mm]{\scriptsize reads / writes} (db);

    % --- ontology ------------------------------------------------------
    \node[store] (onto) at (-5.4,-4.2) {Compiled ontology\\{\scriptsize \texttt{tree\_data/*.json}}};
    \node[off] (build) at (-5.4,-2.1) {Ontology compiler\\{\scriptsize \texttt{build\_from\_ontology.py}}};
    \node[off] (src) at (-5.4,0) {Ontology source\\{\scriptsize skills, edges, topics}};
    \draw[arr] (src) -- (build);
    \draw[arr] (build) -- (onto);
    \draw[arr] (onto) -- node[above]{\scriptsize loaded at boot} (engine);

    % --- content pipeline ----------------------------------------------
    \node[off] (corpus) at (5.4,0) {Source corpus\\{\scriptsize 6 books, 4{,}474 records}};
    \node[off] (rag) at (5.4,-2.1) {RAG retrieval\\{\scriptsize embeddings $+$ top-$k$}};
    \node[off] (gen) at (5.4,-4.2) {Generate $+$ verify\\{\scriptsize two-model pipeline}};
    \draw[arr] (corpus) -- (rag);
    \draw[arr] (rag) -- (gen);
    \draw[arr] (gen) -- node[below, align=center]{\scriptsize validated\\[-1mm]\scriptsize items} (db);
    \draw[arr] (onto.south) |- (-5.4,-7.2) -- node[below]{\scriptsize skill $\times$ Bloom specifications} (5.4,-7.2) -| (gen.south);

    \begin{scope}[on background layer]
      \node[draw, dotted, rounded corners, fit=(src)(build)(onto), inner sep=3mm,
            label={[font=\scriptsize\itshape]below:offline build}] {};
      \node[draw, dotted, rounded corners, fit=(corpus)(rag)(gen), inner sep=3mm,
            label={[font=\scriptsize\itshape]below:offline pipeline}] {};
    \end{scope}
  \end{tikzpicture}
  \caption{System architecture. Solid boxes are the online serving path;
    dashed boxes are offline processes that run before serving and write their
    artefacts into the compiled ontology and the question bank.}
  \label{fig:architecture}
\end{figure}

\begin{table}[htbp]
  \centering
  \caption{Subsystem responsibilities}
  \label{tab:modules}
  \begin{tabularx}{\textwidth}{>{\hsize=0.55\hsize}Y>{\hsize=0.7\hsize}Y>{\hsize=1.75\hsize}Y}
    \toprule
    Subsystem & Location & Responsibility \\
    \midrule
    Knowledge ontology & \path{Ontology/}, \path{Backend/tree_data/} & Defines skills, their Bloom levels, their topics, and the prerequisite edges between them. Compiled from source into the served JSON artefacts by a build script that enforces acyclicity and applies transitive reduction. \\
    Inference engine & \path{Backend/bkt.py}, \path{bloom_taxonomy.py}, \path{mastery_updater.py}, \path{skill.py}, \path{skill_tree.py} & Maintains and updates per-skill mastery; owns the graph, its cycle check and its topological ordering. \\
    Session logic & \path{Backend/diagnostic.py}, session state in \path{server.py} & Runs a diagnostic or practice session: which skill to test next, which Bloom level, which question, when to stop. \\
    Serving API & \path{Backend/server.py} & Exposes catalogue, session and mastery endpoints; composes engine and database. \\
    Content pipeline & \path{data-gen/} & Parses the corpus, builds a retrieval index, generates and verifies questions, loads them to the database. \\
    Client & \path{Frontend/} & Renders the catalogue, question flow and mastery views; handles bilingual display and mathematical typesetting. \\
    \bottomrule
  \end{tabularx}
\end{table}

\paragraph{Where standards shaped the partition.} Two boundaries in
\Cref{fig:architecture} exist for reasons from \Cref{sec:standards} rather than
for reasons of convenience. First, the privileged database key lives only in the
API process: the client never holds a credential that can write to the database,
which is why there is no direct client-to-database arrow even though the
database vendor supports one. Second, the corpus sits strictly inside the
offline pipeline and is never served: generated items reach the database, but
the copyrighted source text never leaves the build environment.

\paragraph{Safety considerations.} The system carries no physical-safety risk,
but it does carry a decision-quality risk: a learner may reallocate study time
during a high-stakes year on the strength of a wrong inference. The design
response is conservatism in the gating policy --- the threshold for declaring a
prerequisite mastered is deliberately high (0.95), so the system prefers to keep
a learner practising a foundation slightly too long over releasing them from it
too early.

\section{Design Refinement and Simulation}
\label{sec:detailed-design}

\subsection{The knowledge-tracing model}
\label{sec:bkt-model}

Each skill's mastery is a latent binary variable; the engine maintains
$P(L_n)$, the probability that the learner has learned the skill after $n$
opportunities. Four parameters govern its evolution: the prior $P(L_0)$, the
learning rate $P(T)$, the slip probability $P(S)$ and the guess probability
$P(G)$ \cite{corbett1995knowledge}. Each response triggers a two-step update.

Step 1, the Bayesian posterior given the observed response:
\begin{align}
  P(L_n \mid \text{correct}) &=
    \frac{P(L_n)\,\bigl(1-P(S)\bigr)}
         {P(L_n)\,\bigl(1-P(S)\bigr) + \bigl(1-P(L_n)\bigr)\,P(G)}
  \label{eq:posterior-correct}\\[2mm]
  P(L_n \mid \text{incorrect}) &=
    \frac{P(L_n)\,P(S)}
         {P(L_n)\,P(S) + \bigl(1-P(L_n)\bigr)\,\bigl(1-P(G)\bigr)}
  \label{eq:posterior-incorrect}
\end{align}

Step 2, the learning transition --- the chance that the opportunity itself
taught the skill:
\begin{equation}
  P(L_{n+1}) = P(L_n \mid \text{evidence})
             + \bigl(1 - P(L_n \mid \text{evidence})\bigr)\,P(T)
  \label{eq:transition}
\end{equation}

Mastery is persisted as a percentage on $[0,100]$ and related to the
probability form by $P(L) = \text{mastery}/100$. The implementation is
deliberately \emph{stateless} --- it takes a probability and an observation and
returns a probability, holding no learner state --- so that the entire model can
be replaced by an individualised or neural variant
\cite{yudelson2013individualized,piech2015deep} without altering a single call
site. That is the concrete design provision for the successor identified in
\Cref{sec:alternatives}.

\subsection{Conditioning the noise parameters on cognitive level}

The refinement that distinguishes this engine from textbook BKT is that $P(G)$
and $P(S)$ are not constants but functions of the question's Bloom level. A
four-option recall item is guessable; an analysis item is not. Conversely, a
learner who knows an analysis-level procedure has more opportunities to slip
while executing it. \Cref{tab:bloom-params} gives the values used.

\begin{table}[htbp]
  \centering
  \caption{Bloom-conditioned guess and slip parameters, and the mastery band each level occupies}
  \label{tab:bloom-params}
  \begin{tabular}{lccc}
    \toprule
    Bloom level & $P(G)$ & $P(S)$ & Mastery band \\
    \midrule
    Remember   & 0.30 & 0.05 & 0.00 -- 16.67 \\
    Understand & 0.25 & 0.08 & 16.67 -- 33.33 \\
    Apply      & 0.18 & 0.12 & 33.33 -- 50.00 \\
    Analyze    & 0.12 & 0.18 & 50.00 -- 66.67 \\
    Evaluate   & 0.08 & 0.22 & 66.67 -- 83.33 \\
    Create     & 0.05 & 0.25 & 83.33 -- 100.00 \\
    \bottomrule
  \end{tabular}
\end{table}

The 0--100 mastery scale is partitioned into six equal bands of
$100/6 \approx 16.67$ points, one per level, so that a mastery value maps
directly to a cognitive level. This is what makes the selection rule expressible:
the next question is requested one band above the learner's current band, which
is the zone-of-proximal-development rule \cite{vygotsky1978mind} rendered
arithmetically.

\subsection{The three-phase mastery update policy}
\label{sec:mastery-policy}

Every response runs three phases in order.

\textbf{Phase 1 --- Bloom-conditioned BKT.} Look up $P(G), P(S)$ for the
question's Bloom level, compute the effective transition rate (Phase 3), and
apply \Cref{eq:posterior-correct,eq:posterior-incorrect,eq:transition} to the
answered skill.

\textbf{Phase 2 --- Ancestor pull-up.} If the response was correct, traverse the
skill's ancestors breadth-first and raise each to at least the answered skill's
new posterior:
\begin{equation}
  P(a) \leftarrow \max\bigl(P(a),\, P(s)\bigr)
  \qquad \forall a \in \operatorname{ancestors}(s)
  \label{eq:pullup}
\end{equation}
The justification is evidential: correctly executing a skill is evidence that
its prerequisites are held, because the skill is not executable without them.
The asymmetry is deliberate --- an \emph{incorrect} response propagates nothing
downward to prerequisites, because failure is ambiguous. It may indicate a
missing prerequisite, but it may equally indicate a slip at the top level, and
penalising every ancestor for one wrong answer would make the estimate
unstable. Locating the responsible prerequisite is instead the job of the
spillover rules, which use option-level evidence rather than a blanket penalty.

\textbf{Phase 3 --- Conjunctive transition gating.} A skill's learning rate is
not constant. For a skill $s$ with direct prerequisites
$\operatorname{parents}(s)$:
\begin{equation}
  P(T)_s =
  \begin{cases}
    \tau_{\text{base}} & \text{if } \min_{p \in \operatorname{parents}(s)} P(L_p) \ge \theta \\[1mm]
    \tau_{\text{locked}} & \text{otherwise}
  \end{cases}
  \label{eq:gating}
\end{equation}
with $\theta = 0.95$, $\tau_{\text{base}} = 0.10$ in normal practice,
$0.20$ during a diagnostic (where the aim is to move the estimate quickly), and
$\tau_{\text{locked}} = 0.01$. The gate is \emph{conjunctive} --- every
prerequisite must clear the threshold, not merely the average or the best one.
This is the formal expression of the project's central claim: a learner cannot
be credited with advancing on a skill while a foundation beneath it is missing.
Its correctness depends entirely on $\operatorname{parents}(s)$ being the
\emph{direct} prerequisites, which is why the ontology build applies transitive
reduction (\Cref{sec:ontology-rebuild}); a redundant shortcut edge would let a
skill unlock while its true intermediate prerequisite was still unmastered.

\subsection{Adaptive question selection}

A diagnostic of tractable length cannot test 1{,}688 skills --- a section's
worth is already far more than a learner will sit through --- so it must choose
skills whose outcome is maximally informative about the rest of the graph. The engine
scores each untested skill by a balanced-volume heuristic:
\begin{equation}
  \operatorname{score}(s) =
  \frac{u_{\uparrow}(s) + u_{\downarrow}(s) + 1}
       {1 + \lvert u_{\uparrow}(s) - u_{\downarrow}(s) \rvert}
  \label{eq:selection}
\end{equation}
where $u_{\uparrow}$ and $u_{\downarrow}$ are the counts of untested ancestors
and untested descendants. The numerator rewards skills that touch many unknowns;
the denominator rewards \emph{balance}. A skill with 20 unknown descendants and
no unknown ancestors is a root whose failure tells us little about what lies
above it; a skill sitting between two large unknown regions bisects the graph,
and its outcome constrains both directions. This is the knowledge-space insight
\cite{doignon1985spaces} applied as a greedy heuristic.

Priors are initialised by depth: root skills start at 0.70 mastery and the
deepest skills at 0.20, interpolated linearly between, encoding the assumption
that foundational material is more likely already held.

Having chosen a skill and a target Bloom level, the engine searches for an item
through a fallback ladder, because the ideal item frequently does not exist in
a partially generated bank: (i) same skill, nearby Bloom level, matching topic;
(ii) same skill, nearby Bloom level, any topic; (iii) same topic, graph-nearby
skills ordered by distance, preferring the lowest-mastery candidate. ``Nearby
Bloom'' means the offsets $[0,+1,-1,+2,-2,+3,-3]$ tried in that order. If no
tier yields an item, the skill is marked tested and the session advances --- the
design decision that a missing question must never stall a session (NFR3).

\subsection{Prerequisite spillover}

During topic practice, two rules detect that failure is caused by something
upstream rather than by the topic itself. Rule A fires when at least three of
the last five wrong answers implicate the same missing prerequisite, evidence
that comes from the \texttt{option\_missing\_prerequisites} table --- each
distractor is annotated with the skill whose absence would lead a learner to
choose it. Rule B fires when rolling accuracy over the last eight attempts falls
below 40\% \emph{and} some candidate prerequisite has mastery below 0.60; the
weakest such prerequisite is chosen. Either rule injects two prerequisite-focused
questions before returning to the topic traversal. Rule A is precise but needs
distractor annotations; Rule B is a coarser safety net that works without them.

\subsection{The skill graph, and its rebuild}
\label{sec:ontology-rebuild}

\begin{figure}[htbp]
  \centering
  \begin{tikzpicture}[
      font=\scriptsize,
      skill/.style={draw, rounded corners=2pt, minimum width=30mm,
                    minimum height=7mm, align=center, fill=gray!8},
      arr/.style={-{Stealth[length=1.8mm]}}
    ]
    \node[skill] (alg)  at (0,0)    {Simplify an algebraic expression\\{\tiny Apply}};
    \node[skill] (diff) at (-3.6,-1.8) {Differentiate a polynomial\\{\tiny Apply}};
    \node[skill] (rec)  at (0.2,-1.8)  {Recognise when parts applies\\{\tiny Understand}};
    \node[skill] (pick) at (0.2,-3.6)  {Choose $u$ and $dv$\\{\tiny Apply}};
    \node[skill] (ibp)  at (-1.7,-5.4) {Evaluate an integral by parts\\{\tiny Apply}};

    \draw[arr] (alg) -- (rec);
    \draw[arr] (rec) -- (pick);
    \draw[arr] (pick) -- (ibp);
    \draw[arr] (diff) -- (ibp);
    \draw[arr] (alg) to[out=180,in=90] (diff);
  \end{tikzpicture}
  \caption{A fragment of the skill graph, illustrating the granularity target.
    ``Integration by parts'' is not one node but a chain: recognising
    applicability, choosing the decomposition, and executing it --- each
    separately testable, and each with its own prerequisites, so that a failure
    can be localised to the step that actually broke.}
  \label{fig:skill-fragment}
\end{figure}

The first ontology was derived from a sample of the corpus and yielded 430
skills across 66 topics with 471 edges. In use it proved too coarse: a topic
held on average fewer than four skills, so a failure inside a topic could not be
localised further, and the conjunctive gate had too few edges to act on. The
graph was therefore rebuilt from the complete corpus --- all 4{,}474 parsed
records of all six books --- at a granularity where a node names one action, and
the rebuild is now the catalogue the application serves.
\Cref{tab:ontology-compare} compares the two. The two edge counts given there
are both real and differ for a reason worth stating: the ontology source asserts
1{,}743 prerequisite relations, and the build's transitive reduction removes the
60 that are implied by a longer path, leaving 1{,}683 direct edges in the
compiled catalogue. It is the compiled figure that \Cref{eq:gating} acts on.

\begin{table}[htbp]
  \centering
  \caption{Legacy ontology versus full-corpus rebuild}
  \label{tab:ontology-compare}
  \begin{tabular}{lrr}
    \toprule
    Property & Legacy (superseded) & Rebuild (served) \\
    \midrule
    Skills & 430 & 1{,}688 \\
    Prerequisite edges (source) & 471 & 1{,}743 \\
    Prerequisite edges (compiled) & 471 & 1{,}683 \\
    Topics covered & 66 & 183 of 183 \\
    Corpus coverage & sample & all 4{,}474 records \\
    Maximum graph depth & 7 & 9 \\
    Isolated skills & --- & 0 \\
    Unresolved prerequisite references & --- & 0 \\
    Connected components (largest) & --- & 87 (482) \\
    \bottomrule
  \end{tabular}
\end{table}

\begin{table}[htbp]
  \centering
  \caption{Rebuilt ontology by subject}
  \label{tab:ontology-subject}
  \begin{tabular}{lrrrr}
    \toprule
    Subject & Skills & Edges & Edges per skill & Topics \\
    \midrule
    Chemistry   & 886 & 868 & 0.97 & 59 \\
    Mathematics & 446 & 499 & 1.10 & 52 \\
    Physics     & 356 & 376 & 1.06 & 72 \\
    \midrule
    Total       & 1{,}688 & 1{,}743 & 1.03 & 183 \\
    \bottomrule
  \end{tabular}
\end{table}

The build is reproducible rather than hand-maintained: a compiler reads the
ontology source, applies an editorial layer (topic aliases and display labels),
checks the graph for cycles, applies transitive reduction, and emits the JSON
artefacts the server loads. Transitive reduction is not cosmetic --- as
\Cref{eq:gating} shows, a shortcut edge would let a skill unlock while its true
intermediate prerequisite was unmastered.

\subsection{Data model}

\begin{table}[htbp]
  \centering
  \caption{Database schema}
  \label{tab:schema}
  \begin{tabularx}{\textwidth}{>{\hsize=0.65\hsize}Y>{\hsize=1.35\hsize}Y}
    \toprule
    Table & Contents and constraints \\
    \midrule
    \texttt{users} & Learner identity: surrogate UUID, unique username. \\
    \texttt{user\_skill} & Per-learner, per-skill mastery on $[0,100]$; composite primary key; cascades on learner or skill deletion. \\
    \texttt{skills} & Skill identifier and description; description is NOT NULL. \\
    \texttt{questions} & Stem, Bloom level, and topic. \texttt{topic} is a foreign key to \texttt{ontology\_topics}, so it stores a topic \emph{code}, never a display label. \\
    \texttt{question\_options} & Four options per question, label constrained to A--D, correctness flag, and a per-option explanation. \\
    \path{option_missing_prerequisites} & Maps a distractor to the skill whose absence explains choosing it --- the evidence source for spillover Rule A. \\
    \texttt{ontology\_topics}, \texttt{ontology\_skill\_topics}, \texttt{ontology\_skill\_edges} & The ontology mirrored into the database for referential integrity of questions and skills. \\
    \bottomrule
  \end{tabularx}
\end{table}

Two deliberate omissions are worth stating. There is no \texttt{subject} column
anywhere: subject is a property of the ontology, derived at read time and
plumbed through the API, so that the database never holds a second copy of a
fact the ontology owns. And the transition probability $P(T)_s$ is not
persisted: it is recomputed from current parent mastery on every request, so it
can never fall out of sync with the values it is derived from.

\subsection{API surface}

The API is a flat, session-oriented REST surface of 14 endpoints
(\Cref{tab:api}). The two session families are deliberately symmetric ---
\texttt{start}, \texttt{status}, \texttt{next}, \texttt{answer} --- so the client
can drive either flow with one component.

\begin{table}[htbp]
  \centering
  \caption{API endpoints}
  \label{tab:api}
  \begin{tabularx}{\textwidth}{>{\hsize=1.05\hsize}Y>{\hsize=0.95\hsize}Y}
    \toprule
    Endpoint & Purpose \\
    \midrule
    \texttt{GET /catalog} & Courses, sections and topics from the compiled ontology \\
    \texttt{GET /users/\{name\}/}, \texttt{POST /users/} & Look up and create a learner \\
    \texttt{GET .../sections/\{id\}/state} & Diagnostic gate and lock state for a section \\
    \texttt{GET .../sections/\{id\}/mastery} & Mastery graph and table for visualisation \\
    \texttt{POST .../diagnostic/retake} & Reset section mastery and clear its sessions \\
    \texttt{POST /diagnostic/start} & Begin a diagnostic session \\
    \texttt{GET /diagnostic/\{id\}/status}, \texttt{/next} & Session progress; fetch next selected question \\
    \texttt{POST /diagnostic/\{id\}/answer} & Submit a response; triggers the three-phase update \\
    \texttt{POST /topic-practice/start} & Begin practice on one topic \\
    \texttt{GET /topic-practice/\{id\}/status}, \texttt{/next} & Progress; next question, including injected spillover items \\
    \texttt{POST /topic-practice/\{id\}/answer} & Submit a response; evaluates spillover rules \\
    \bottomrule
  \end{tabularx}
\end{table}

\subsection{Content generation pipeline}
\label{sec:generation-pipeline}

The pipeline turns the corpus into items aligned to the ontology. Source books
are parsed into structured records; a multilingual embedding model
\cite{reimers2019sentence} indexes them; for each (skill, Bloom level)
specification the top-$k$ relevant records are retrieved and supplied to a
generator model as grounding \cite{lewis2020retrieval}.

The design decision that matters most is the second stage. A generated item is
not accepted on the generator's own word: a \emph{different, cheaper} model
independently checks that the item is on-topic for the claimed skill, and
\emph{blind re-solves} it --- it is shown the question and options but not the
marked answer --- and the item is accepted only if its independent answer agrees
with the marked one. Letting a model grade its own output biases towards false
confidence; an independent judge with no sight of the expected answer is a real
check. The cost asymmetry is deliberate: generation uses the stronger model
because it is the harder task, verification uses the cheaper one because
agreement checking is easier than authoring, and verification calls are batched
across several specifications to reduce call count.

\subsection{Design calculations}

\textbf{Generation space.} With one item per (skill, Bloom level) pair, the
legacy ontology implies $430 \times 6 = 2{,}580$ specifications; the rebuilt
ontology implies $1{,}688 \times 6 = 10{,}128$. At the observed generation
pattern this is the dominant cost driver in the project and the reason adoption
of the rebuilt ontology is gated on a budgeted generation run
(\Cref{sec:budget}).

\textbf{Storage.} A question with four options and explanations occupies on the
order of a few kilobytes. A complete bank at rebuild scale is therefore of the
order of tens of megabytes --- immaterial against any database tier. Learner
data is smaller still: one row per learner per skill, so a cohort of
$10{,}000$ learners against 1{,}688 skills is under $2\times10^7$ rows of two
identifiers and a float, which is well within commodity capacity.

\textbf{Inference cost.} One response costs one constant-time Bayesian update
plus an ancestor traversal bounded by graph depth (9) and a successor refresh
over direct children. The engine is not the scaling bottleneck; database round
trips during question search are.

\section{Application of Tools}
\label{sec:tools}

% Rendered as a longtable (not a float) so it flows across pages at readable
% size instead of being squeezed onto one rotated page.
\setlength{\LTcapwidth}{\textwidth}
{\small\setstretch{1.1}
\begin{longtable}{@{}p{30mm}p{52mm}p{\dimexpr\textwidth-30mm-52mm-4\tabcolsep\relax}@{}}
  \caption{Tool selection, justification, and relevant limitations}
  \label{tab:tools}\\
    \toprule
    Tool & Why selected & Relevant limitations \\
    \midrule
  \endfirsthead
    \multicolumn{3}{@{}l}{\small\itshape \tablename~\thetable{} --- continued from previous page}\\[2pt]
    \toprule
    Tool & Why selected & Relevant limitations \\
    \midrule
  \endhead
    \midrule
    \multicolumn{3}{r@{}}{\small\itshape continued on next page}\\
  \endfoot
    \bottomrule
  \endlastfoot
    Python 3 & The engine is mathematical and graph-oriented; Python keeps \Cref{eq:posterior-correct}--\Cref{eq:gating} readable enough to audit against the policy document. & Slower than a compiled language --- immaterial here, since the update is constant-time. \\
    FastAPI $+$ Uvicorn & Typed request/response models via Pydantic catch malformed payloads at the boundary; automatic OpenAPI documentation made the API self-describing for the frontend developers. & The asynchronous model gives no benefit as used, since the engine is synchronous and CPU-light. \\
    Supabase (PostgreSQL) & A managed relational store with a usable free tier. Relational integrity is not optional here --- foreign keys are what stop a question referencing a skill that no longer exists. & The free tier pauses idle projects; the client library adds a network round trip where raw SQL would allow one query. \\
    React 18 $+$ Vite & Component model suits a question-at-a-time flow with shared state; Vite's fast rebuilds mattered on laptop hardware. & Single-page delivery means a large first-load bundle, worsened by mathematical typesetting fonts. \\
    KaTeX & Renders \LaTeX{} mathematics synchronously and quickly, and outputs accessible markup rather than images \cite{wcag21}. & Supports a subset of \LaTeX{}; unsupported macros in generated content must be caught before serving. \\
    Tailwind CSS & Utility classes kept styling consistent across developers without a design system. & Verbose markup; class strings are hard to review in diffs. \\
    Gemini API & After local generation proved infeasible on CPU (\Cref{sec:feasibility}), a hosted model supplied the throughput; free-tier keys could be pooled to raise it further. & Rate limits shape the pipeline's architecture; a hosted model is a dependency the project does not control, and model updates can silently change output quality. \\
    Sentence-transformers & Multilingual embeddings are required for a Bangla corpus; English-only models retrieve poorly here \cite{reimers2019sentence}. & A large download and CPU-bound indexing; embedding quality on Bangla mathematical text is materially worse than on English prose. \\
    Git and GitHub & Version control and the integration point for five developers working on separable subsystems. & --- \\
\end{longtable}}

\section{Implementation}
\label{sec:implementation}

\textbf{Engine.} The model was implemented from the equations rather than from a
library, which was the right call for auditability: the policy document and the
code state the same constants, and a reviewer can check one against the other.
The graph enforces its own invariant --- adding a prerequisite edge provisionally
inserts it, runs a three-colour depth-first search over the whole graph, and
rolls the edge back if it would close a cycle. The graph cannot enter an invalid
state through the public interface.

\textbf{Ontology.} Generated artefacts are never hand-edited. The source is
compiled by script, and the editorial decisions that cannot be derived from data
--- topic aliases, display labels, course and section layout --- live in a single
configuration file. Rebuilding is a one-command operation, which is what made a
late full rebuild affordable.

\textbf{Content pipeline.} Two failures shaped the implementation more than any
design intent. The first: source books are machine-extracted, and one corrupt
record in the Chemistry corpus --- a string containing a \LaTeX{} macro repeated
tens of thousands of times and never terminated --- desynchronised the parser and
destroyed 334 neighbouring records. The parser was rewritten to resynchronise
by decoding one object at a time, so a bad record now costs only itself. The
second: language models emit \LaTeX{} with single backslashes, and a strict JSON
parser silently consumes several of those escapes as control characters rather
than raising an error, corrupting text invisibly. This affected a measured
13.4\% of one early batch. The fix is an escape-repair pass applied before every
parse, in both the ingest and generation paths. Both failures are documented in
the repository so that a future contributor does not rediscover them.

\textbf{Frontend.} The mathematics rendering component is the non-obvious piece:
it splits mixed Bangla-and-\LaTeX{} strings on delimiters, renders the
mathematical segments through KaTeX, renders the rest as text, and normalises
newline sequences that arrive literally from generated content.

\textbf{Process.} Work was carried out in a shared Git repository with
subsystem-aligned branches merged into a main line --- 93 commits across roughly
five and a half months from five contributors. In place of a formal CI service,
the project relies on a credential-free test harness that any developer can run
in seconds before merging (\Cref{sec:evaluation}), plus a chain of handover
documents that record the state and rationale of each subsystem at the point it
changed hands. Those documents are themselves an engineering artefact: they
exist because the project spanned several months and five people, and knowledge
that lived only in one contributor's head was repeatedly found to be the
project's real bottleneck.

\section{Deployment}
\label{sec:deployment}

\textbf{Target environment.} The system deploys as two processes and one managed
service: a Python ASGI application serving the API, a static bundle for the
client, and the hosted PostgreSQL instance. A container definition is provided
for the backend. Configuration is by environment variables --- the database URL
and service key --- supplied through a file that is excluded from version
control; the client takes the API base URL at build time and defaults to the
local backend.

\textbf{Release process.} The ontology is compiled and its report inspected;
generated questions are validated and loaded with a dry run first; the backend
and client are then started. The deliberate ordering reflects the dependency
chain in \Cref{fig:architecture}: content cannot be loaded against an ontology
that has not been compiled.

\textbf{Operation and maintenance.} The prototype is operated by the project
team. Monitoring is limited to structured logging of question-search behaviour
--- each search prints which tier matched and how long the database query took,
which is what made the topic-code defect in \Cref{sec:revisions} visible.
Backup is delegated to the managed database service.

\textbf{Resources consumed.} One small application process (well under a
gigabyte of memory, since the loaded graph is a few thousand nodes), one static
file bundle, and a database holding tens of megabytes. Offline, the pipeline
consumes CPU for embedding the corpus and external API calls for generation ---
the only materially resource-hungry part of the system, and the subject of the
environmental assessment in \Cref{sec:sustainability}.
